% 01_intro.tex  --  Introduction

\section{Introduction}
\label{sec:intro}

ML engineers and data scientists who train gradient-boosted models on tabular
data face a recurrent resource tension. A single retraining run on a dataset of
hundreds of thousands of rows can consume substantial compute, and repeated
runs for hyperparameter search, model versioning, or data-sharing workflows
multiply that cost many times over. SLD is cheap for a different reason: it
costs one histogram pass under a fixed prediction state, before retraining on
each candidate surrogate. Condensation is not itself a privacy or compliance
guarantee, but a small surrogate can be useful when combined with appropriate
governance or privacy mechanisms. The natural remedy, replacing the
full dataset with a small, high-fidelity surrogate, is already practiced
informally through random subsampling and stratified sampling. What is missing
is a principled way to answer the question that matters most \emph{before}
committing to retraining on a surrogate: does this condensed set faithfully
preserve the structure that the booster would have learned from the original
data?

The dataset condensation and coreset-selection literature has produced a rich
supply of methods to construct small surrogates~\cite{wang2018datasetdistillation,
zhao2021dcgm,welling2009herding,mirzasoleiman2020craig}. For tabular data
specifically, several methods have recently been adapted or purpose-built for
non-image settings~\cite{kang2025tabdistill,xu2026c2tc}. Yet the evaluation
paradigm is almost universally the same: train a fresh model on the condensed
set, measure downstream predictive performance, and compare. This protocol has two practical
shortcomings. First, it requires precisely the retraining one hoped to avoid.
Second, it collapses a mechanistic question, namely \emph{why} a condensed set
works or fails, into a scalar outcome that gives no guidance for debugging or
selection.

Gradient-boosted decision trees~\cite{friedman2001gbm,chen2016xgboost,
ke2017lightgbm} offer a distinctive opportunity to break this cycle. A GBDT
tree structure is determined by split-selection decisions: at each node, the
algorithm chooses the feature and threshold that maximize the regularized gain
criterion (\cref{sec:diagnostic}). Predictive behaviour also depends on leaf
values, learning rate, boosting rounds, missing-value handling, and
regularization. The gain criterion is histogram-sufficient once first- and
second-order loss statistics are fixed; those statistics may come from a root
base score or from a reference booster. Comparing the histogram landscapes of
the full and condensed datasets is therefore a no-candidate-retraining screening
diagnostic, not a claim that no model state is ever used.

We introduce the \emph{Split-Landscape Discrepancy} (SLD), which formalizes
this comparison as an $\ell_p$ distance between the split-gain landscapes of a
full dataset and its surrogate (\cref{sec:diagnostic}). Using SLD as a
measurement instrument, we conduct a systematic comparative evaluation of 17
condensation methods, spanning classical coresets, gradient-based selectors,
distribution-matching distillation, and gain-path distillation. The 17-method
audit uses 26 datasets, two budgets, and three seeds; the fuller five-budget,
five-seed core grid is restricted to five methods. The analysis is a
measurement study; we do not propose a new condensation method and we do not
claim a new state of the art.

Our analysis surfaces three empirical findings of practical consequence.
At the method-mean level, split-regret strongly separates method families by
downstream AUROC in the small-to-mid-size benchmark, enabling rejection
filtering before retraining on every candidate.
A provably coverage-optimal selector is simultaneously among the weakest methods
by predictive performance, a demonstration that coverage and structural fidelity
are independent objectives. At deployment scale, random sampling becomes
difficult to beat, and SLD is better read as a rejection signal than as a
single-best selector.

\smallskip
\noindent\textbf{Contributions.}
\begin{itemize}[leftmargin=1.2em,itemsep=2pt]
  \item \textbf{The SLD diagnostic.} We introduce the Split-Landscape
    Discrepancy, a no-candidate-retraining metric that quantifies how
    faithfully a condensed dataset reproduces the GBDT split-gain landscape of
    the full data under a fixed prediction state (\cref{sec:diagnostic}).

  \item \textbf{Systematic comparative analysis.} We present a 17-method
    diagnostic audit for GBDT-on-tabular settings and a separate five-method
    core grid over 26 datasets, five budgets, and five seeds, measuring both
    downstream predictive performance and landscape fidelity
    (\cref{sec:setup,sec:results}).

  \item \textbf{Three empirically grounded insights.} A single landscape score
    is method-level rank-correlated with AUROC at Spearman $\rho = -0.87$ in
    the small-to-mid-size diagnostic grid, making it useful for screening and
    family-level ordering rather than top-1 selection. Within-dataset trends are
    generally negative after controlling for dataset scale. Coverage and predictive
    fidelity are decoupled: a $(1-1/e)$-optimal coverage selector is among the
    weakest downstream performers, explained by a reference-chain decomposition
    into structure and leaf-estimate gaps. Structure is robust to perturbation
    below a margin-linked threshold; the dose is a fraction of the split margin,
    not a fraction of intentionally corrupted splits (\cref{sec:results}).

  \item \textbf{Scope boundaries.} We document where the lens does not extend:
    regression tasks (structure-dominated failure), hyperparameter-rank
    transfer (near-zero Jaccard), and no-retraining selection of the single
    best method (top-3 rate $0.44$), measured under the same protocol
    (\cref{sec:limits}).

  \item \textbf{Released artifacts.} The benchmark protocol, SLD
    implementation, processed tables, seeds, command lines, and plotting scripts
    are released at
    \url{https://github.com/VIVAN-RANGRA/sld-gbdt-condensation}.
\end{itemize}
